\chapter{Aprendizado profundo para séries temporais}
\label{cap:profundo}

\begin{objetivos}
Este capítulo é um panorama de referência, não um curso de aprendizado profundo: ele
apresenta apenas o que é necessário para acompanhar as arquiteturas usadas no restante
do livro. Ao final, você será capaz de:
\begin{itemize}[leftmargin=*,itemsep=0.2em]
  \item descrever a estrutura de um perceptron de múltiplas camadas e explicar por que
        ele é inadequado a séries temporais longas;
  \item explicar, no nível das equações de estado, o que as células LSTM e GRU fazem e
        que problema de gradiente elas atenuam;
  \item aplicar convoluções unidimensionais a sinais e calcular o campo receptivo de uma
        pilha convolucional;
  \item descrever a estrutura de um autoencoder, interpretar seu espaço latente e usar o
        erro de reconstrução como escore de anomalia;
  \item distinguir autoencoder, autoencoder variacional e rede de agrupamento profundo
        quanto ao que cada um otimiza;
  \item descrever a arquitetura U-Net e justificar por que suas conexões de atalho
        importam para segmentação;
  \item enunciar a equação da atenção e situar o Transformer no contexto deste livro ---
        nenhum modelo aqui treinado é um Transformer, e a arquitetura não é desenvolvida
        em detalhe.
\end{itemize}
\end{objetivos}

\section{Da regressão à rede profunda}

\subsection{O perceptron de múltiplas camadas}

Uma camada densa aplica uma transformação afim seguida de não linearidade:
\begin{equation}
  \vect{h}^{(\ell)} = \sigma\big( \mat{W}^{(\ell)} \vect{h}^{(\ell-1)} + \vect{b}^{(\ell)} \big),
  \label{eq:camada-densa}
\end{equation}
com $\vect{h}^{(0)} = \vect{x}$. Empilhando $L$ camadas obtém-se o \termo{perceptron de
múltiplas camadas}{multilayer perceptron} (MLP). As não linearidades usadas neste livro
são a ReLU, $\sigma(z) = \max(0, z)$, nas camadas ocultas, e a sigmoide,
$\sigma(z) = 1/(1 + e^{-z})$, nas saídas binárias.

Aplicar um MLP a uma série temporal exige achatá-la em um vetor de dimensão $T \cdot S$.
Para $T = 200$ e $S = 4$, isso resulta em \num{800} entradas --- exatamente a
representação ``bruta'' usada como linha de base no \cref{cap:mahalanobis}. O
procedimento funciona, mas descarta duas propriedades essenciais do sinal: a
\emph{ordem} (permutar os instantes produziria a mesma representação para uma rede sem
estrutura) e a \emph{invariância a translação} (o mesmo padrão em posição diferente
ativa pesos diferentes).

\begin{destaque}[Por que a linha de base bruta ainda aparece]
O \cref{cap:mahalanobis} reporta que classificadores binários sobre a representação
bruta de \num{800} dimensões atingem $0{,}691 \pm 0{,}162$ de acurácia, contra
$0{,}755 \pm 0{,}171$ sobre a representação latente de um classificador multiclasse. A
diferença é real, mas modesta --- e o fato de a representação bruta chegar tão perto é
informativo: as assinaturas dessas anomalias são, em boa medida, capturáveis por
estatísticas de primeira ordem.
\end{destaque}

\subsection{Regularização}

Três mecanismos são usados de forma recorrente:

\begin{itemize}[leftmargin=*,itemsep=0.3em]
  \item \textbf{\ing{Dropout}} \citep{srivastava14a}: durante o treinamento, cada
        unidade é zerada com probabilidade $p$. Nas arquiteturas deste livro, $p$ varia
        de $0{,}1$ a $0{,}5$, tipicamente decrescendo em direção à saída.
  \item \textbf{Normalização em lote} (\ing{batch normalization}): normaliza as ativações
        de cada camada pela média e variância do lote, acelerando a convergência e
        atuando como regularizador leve. Presente em todos os blocos da U-Net do
        \cref{cap:segmentacao}.
  \item \textbf{Parada antecipada} (\ing{early stopping}): interrompe o treinamento
        quando a perda de validação deixa de melhorar por um número fixo de épocas.
        O \cref{cap:closedworld} usa paciência de cinco épocas.
\end{itemize}

\section{Redes recorrentes}

\subsection{O problema}

Uma rede recorrente processa a sequência mantendo um estado oculto:
\begin{equation}
  \vect{h}_t = \sigma\big( \mat{W}_{hh}\vect{h}_{t-1} + \mat{W}_{xh}\vect{x}_t + \vect{b} \big).
  \label{eq:rnn}
\end{equation}
A retropropagação através do tempo desdobra essa recorrência em $T$ passos, e o
gradiente em relação a $\vect{h}_1$ envolve o produto de $T$ jacobianas. Se os autovalores
dominantes dessas matrizes forem menores que um, o gradiente decai exponencialmente e a
rede não aprende dependências longas; se forem maiores, ele explode. É o
\termo{problema do gradiente evanescente}{vanishing gradient problem}
\citep{SCHMIDHUBER201585}.

Para $T = 500$, como no \cref{cap:dual}, uma RNN simples é inviável.

\subsection{A célula LSTM}

A \termo{memória de curto prazo longa}{long short-term memory} (LSTM)
\citep{Hochreiter1997} introduz um estado de célula $\vect{c}_t$ que se propaga por
soma, não por produto matricial, e três portas que regulam o fluxo de informação:

\begin{align}
\vect{f}_t &= \sigma\big(\mat{W}_f [\vect{h}_{t-1}, \vect{x}_t] + \vect{b}_f\big)
  && \text{(esquecimento)} \label{eq:lstm-f}\\
\vect{i}_t &= \sigma\big(\mat{W}_i [\vect{h}_{t-1}, \vect{x}_t] + \vect{b}_i\big)
  && \text{(entrada)} \label{eq:lstm-i}\\
\tilde{\vect{c}}_t &= \tanh\big(\mat{W}_c [\vect{h}_{t-1}, \vect{x}_t] + \vect{b}_c\big)
  && \text{(candidato)} \label{eq:lstm-c}\\
\vect{c}_t &= \vect{f}_t \odot \vect{c}_{t-1} + \vect{i}_t \odot \tilde{\vect{c}}_t
  && \text{(estado de célula)} \label{eq:lstm-cell}\\
\vect{o}_t &= \sigma\big(\mat{W}_o [\vect{h}_{t-1}, \vect{x}_t] + \vect{b}_o\big)
  && \text{(saída)} \label{eq:lstm-o}\\
\vect{h}_t &= \vect{o}_t \odot \tanh(\vect{c}_t)
  && \text{(estado oculto)} \label{eq:lstm-h}
\end{align}

onde $\odot$ denota produto elemento a elemento. A \cref{eq:lstm-cell} é o coração do
mecanismo: quando $\vect{f}_t \approx \vect{1}$ e $\vect{i}_t \approx \vect{0}$, o
estado de célula é copiado inalterado, e o gradiente flui sem atenuação.

\begin{destaque}[LSTM em todos os capítulos da Parte II e III]
O \cref{cap:dual} usa uma pilha LSTM(16)--LSTM(4) como codificador e o espelho como
decodificador. O \cref{cap:mundoaberto} usa LSTM(32)--LSTM(8) com camada densa latente
de dimensão 16. O \cref{cap:mahalanobis} usa LSTM(256)--LSTM(128) com latente de
dimensão 128. A tendência de crescimento não é acidental: reflete o aumento da
quantidade de dados e a mudança de objetivo, de reconstrução para representação
discriminativa.
\end{destaque}

\subsection{GRU}

A \termo{unidade recorrente com portas}{gated recurrent unit} funde as portas de
esquecimento e entrada em uma única porta de atualização, eliminando o estado de célula
separado. Tem cerca de um quarto a menos de parâmetros e desempenho comparável na
maioria das tarefas. Não é usada nos experimentos deste livro, mas é a primeira
alternativa a considerar quando o custo computacional importa.

\section{Convoluções unidimensionais}

Uma convolução unidimensional com núcleo de tamanho $\kappa$ aplica
\begin{equation}
  h_t^{(j)} = \sigma\left( \sum_{s=1}^{S} \sum_{u=0}^{\kappa-1}
  w^{(j)}_{s,u}\, x_{t+u}^{(s)} + b^{(j)} \right),
  \label{eq:conv1d}
\end{equation}
compartilhando os pesos $w^{(j)}$ ao longo de todo o eixo temporal. Duas consequências
seguem imediatamente: o número de parâmetros independe de $T$, e a resposta é
equivariante a translações --- o mesmo padrão detectado em posições diferentes produz a
mesma ativação, deslocada.

\begin{definicao}{Campo receptivo}
O campo receptivo de uma unidade é o número de amostras da entrada que influenciam seu
valor. Para uma pilha de $L$ camadas convolucionais de núcleo $\kappa$, sem
subamostragem, o campo receptivo é $L(\kappa - 1) + 1$. Com subamostragem por fator 2
após cada camada, ele cresce aproximadamente como $2^{L}$.
\end{definicao}

\begin{exemplo}{O campo receptivo da U-Net do \cref{cap:segmentacao}}
A arquitetura tem dois blocos de codificação com duas convoluções de núcleo 5 cada,
seguidos de subamostragem por 2, e um gargalo com mais duas convoluções de núcleo 5.
Sem subamostragem, seriam $6 \times 4 + 1 = 25$ amostras. Com as duas subamostragens,
o campo receptivo efetivo no gargalo ultrapassa oitenta amostras, cobrindo boa parte da
janela de duzentas.

O cálculo importa: se o campo receptivo fosse menor que a duração característica do
evento, a rede não poderia, em princípio, reconhecê-lo.
\end{exemplo}

\subsection{Convolução \emph{versus} recorrência}

\begin{center}
\small
\begin{tabular}{p{3.2cm} p{4.7cm} p{4.7cm}}
\toprule
& \textbf{Recorrente (LSTM)} & \textbf{Convolucional 1D} \\
\midrule
Paralelização & Sequencial; lenta em $T$ grande. & Totalmente paralelizável. \\
Dependência longa & Ilimitada em princípio. & Limitada pelo campo receptivo. \\
Parâmetros & Independem de $T$. & Independem de $T$. \\
Adequação & Dinâmica com memória de estado. & Padrões locais e invariantes a
translação. \\
Uso neste livro & \cref{cap:dual}, \cref{cap:mundoaberto},
\cref{cap:mahalanobis}. & \cref{cap:segmentacao}. \\
\bottomrule
\end{tabular}
\end{center}

\section{Autoencoders}

\begin{definicao}{Autoencoder}
Um \termo{autoencoder}{autoencoder} é um par de funções
$g_{\phi}: \R^{T \times S} \to \R^{d}$ (codificador) e
$h_{\psi}: \R^{d} \to \R^{T \times S}$ (decodificador), treinado para minimizar o erro
de reconstrução
\begin{equation}
  \mathcal{L}(\phi, \psi) = \frac{1}{n}\sum_{i=1}^{n}
  \mathcal{D}\big( \mat{X}_i,\; h_{\psi}(g_{\phi}(\mat{X}_i)) \big),
  \label{eq:ae}
\end{equation}
com $d \ll T \cdot S$. O vetor $\vect{z} = g_{\phi}(\mat{X})$ é a
\termo{representação latente}{latent representation}.
\end{definicao}

\begin{figure}[htbp]
\centering
\includegraphics[width=0.85\textwidth]{autoencoder.png}
\caption{Estrutura de um autoencoder: o codificador comprime a entrada até uma
representação latente de baixa dimensão e o decodificador tenta reconstruí-la. A
restrição $d \ll T\cdot S$ força a rede a reter apenas os fatores de variação mais
relevantes.}
\label{fig:profundo-ae}
\end{figure}

A restrição dimensional é o que dá poder ao mecanismo. Como $\vect{z}$ não tem
capacidade para armazenar a entrada integralmente, o treinamento força a retenção dos
fatores de variação mais frequentes. Se o treinamento usa apenas dados normais, esses
fatores são os da normalidade, e a reconstrução de uma entrada anômala é ruim por
construção.

\begin{definicao}{Escore de anomalia por reconstrução}
Dado um autoencoder treinado sobre dados normais,
\begin{equation}
  a(\mat{X}) = \mathcal{D}\big( \mat{X},\; h_{\psi}(g_{\phi}(\mat{X})) \big),
  \label{eq:score-recon}
\end{equation}
e declara-se anomalia quando $a(\mat{X}) > \tau$. O limiar é tipicamente estimado a
partir da distribuição de $a$ sobre dados normais, na forma $\tau = \mu_a + k\sigma_a$.
\end{definicao}

O \cref{cap:dual} usa $\mathcal{D} = \text{SMAPE}$ e $k = 3$. O \cref{cap:mahalanobis}
abandona esse escore em favor da distância no espaço latente, e o motivo será
desenvolvido lá: um autoencoder treinado para reconstruir pode reconstruir bem demais,
inclusive anomalias, quando sua capacidade é suficiente.

\subsection{Autoencoder variacional}

O \termo{autoencoder variacional}{variational autoencoder} (VAE) \citep{VAE,Kingma2014}
substitui o vetor latente determinístico por uma distribuição. O codificador produz
$\vect{\mu}(\mat{X})$ e $\vect{\sigma}(\mat{X})$, amostra-se
$\vect{z} \sim \mathcal{N}(\vect{\mu}, \operatorname{diag}(\vect{\sigma}^2))$ pelo
truque de reparametrização, e acrescenta-se à perda um termo de divergência de
Kullback-Leibler que puxa a posterior aproximada para uma normal padrão:
\begin{equation}
  \mathcal{L}_{\text{VAE}} = \underbrace{\E_{q}\big[ -\log p(\mat{X} \mid \vect{z}) \big]}_{\text{reconstrução}}
  + \beta\, \underbrace{D_{\mathrm{KL}}\big( q(\vect{z}\mid\mat{X}) \,\|\, \mathcal{N}(\vect{0}, \mat{I}) \big)}_{\text{regularização}}.
  \label{eq:vae}
\end{equation}

O efeito é um espaço latente contínuo e suave. A expectativa natural seria que isso
melhorasse o agrupamento subsequente; o \cref{cap:mundoaberto} mostra que não melhorou:
$0{,}727$ contra $0{,}730$ do autoencoder determinístico, diferença sem significância.

\subsection{Rede de agrupamento profundo}

A \termo{rede de agrupamento profundo}{deep clustering network} (DCN) \citep{DCN} otimiza
simultaneamente reconstrução e uma perda de agrupamento no espaço latente:
\begin{equation}
  \mathcal{L}_{\text{DCN}} = \mathcal{L}_{\text{recon}} +
  \lambda \sum_{i} \big\| \vect{z}_i - \vect{c}_{\pi(i)} \big\|^2,
  \label{eq:dcn}
\end{equation}
onde $\vect{c}_{\pi(i)}$ é o centroide atribuído a $\vect{z}_i$. A intenção é produzir
um espaço latente com agrupamentos bem separados. No \cref{cap:mundoaberto}, obteve
$0{,}725$ --- também sem ganho.

\begin{figure}[htbp]
\centering
\includegraphics[width=0.72\textwidth]{DCN_latent.png}
\caption{Projeção t-SNE do espaço latente produzido por uma DCN. Cada cor é uma classe
conhecida; o polígono destaca a região ocupada pelas janelas da falha de ICV, uma classe
que a rede nunca viu. Os pontos novos são absorvidos pelos agrupamentos existentes em vez
de formar um agrupamento próprio.}
\label{fig:profundo-dcn}
\end{figure}

A \cref{fig:profundo-dcn} mostra por que o ganho não aparece --- e o motivo é mais
interessante que a ausência de ganho. A perda de agrupamento da \cref{eq:dcn} pune
explicitamente pontos distantes de qualquer centroide. Diante de uma classe nunca vista,
a rede faz exatamente o que foi instruída a fazer: puxa os pontos novos para dentro dos
agrupamentos que já conhece.

\begin{destaque}[Regularizar o latente pode ser contraproducente]
Autoencoder, VAE e DCN diferem justamente em quanto impõem estrutura ao espaço latente ---
nenhuma, gaussiana e por agrupamentos, respectivamente. Para reconstruir, mais estrutura
ajuda. Para \emph{detectar o desconhecido}, ela atrapalha.

O autoencoder simples reconstrói mal o que nunca viu, e essa incompetência é útil: joga
os pontos novos para as bordas do latente, onde são fáceis de isolar. O VAE os
regulariza para junto dos conhecidos; a DCN os absorve nos agrupamentos. Ambos apagam o
sinal que se queria capturar.

É um caso raro em que a propriedade desejável de um modelo é a sua fragilidade.
\end{destaque}

\begin{destaque}[Um resultado negativo instrutivo]
Três variantes de autoencoder, três resultados indistinguíveis: $0{,}730$, $0{,}727$ e
$0{,}725$. A conclusão do \cref{cap:mundoaberto} é que a escolha da variante não era o
gargalo. O gargalo estava no \emph{objetivo}: todas as três otimizam reconstrução, e
reconstrução não é discriminação.

O \cref{cap:mahalanobis} testa a hipótese oposta --- extrair a representação da camada
penúltima de um classificador multiclasse, treinado com objetivo discriminativo --- e
obtém uma melhora de outra ordem de grandeza: índice de silhueta de $0{,}2299$ contra
$-0{,}0149$, e razão de separação de até 186,5 vezes contra menos de 2,5.
\end{destaque}

\section{A U-Net unidimensional}

A U-Net \citep{Ronneberger2015} foi proposta para segmentação de imagens biomédicas e
adaptada aqui para segmentação de séries temporais. Sua estrutura tem três partes.

\begin{description}[leftmargin=0pt,style=nextline,itemsep=0.4em]

\item[Caminho de contração.] Blocos de convolução seguidos de subamostragem. A cada
bloco, a resolução temporal cai pela metade e o número de canais dobra. Captura contexto
progressivamente mais amplo.

\item[Caminho de expansão.] Blocos de superamostragem seguidos de convolução, que
restauram a resolução original.

\item[Conexões de atalho.] Cada nível do caminho de expansão recebe, por concatenação, o
mapa de ativação do nível correspondente do caminho de contração.

\end{description}

\begin{destaque}[Por que as conexões de atalho são o ponto principal]
Sem elas, a informação de localização precisa seria destruída pelas subamostragens
sucessivas: a rede saberia \emph{que} há uma anomalia na janela, mas não \emph{onde} ela
começa. As conexões de atalho reinjetam, em alta resolução, exatamente essa informação.

Para segmentação temporal de eventos de poço, isso é a diferença entre ``esta janela de
duzentos segundos contém um evento'' e ``o evento começa no segundo 137''. A segunda
resposta é a que o operador precisa.
\end{destaque}

A instanciação concreta, com dois níveis, núcleos de tamanho 5 e 32, 64 e 128 canais, é
detalhada no \cref{cap:segmentacao}.

\section{Atenção e Transformers}

O mecanismo de \termo{atenção}{attention} \citep{Vaswani2017} calcula, para cada
posição, uma combinação ponderada de todas as posições da sequência:
\begin{equation}
  \text{Atenção}(\mat{Q}, \mat{K}, \mat{V}) =
  \operatorname{softmax}\!\left( \frac{\mat{Q}\mat{K}^{\top}}{\sqrt{d_k}} \right) \mat{V},
  \label{eq:atencao}
\end{equation}
onde $\mat{Q}$, $\mat{K}$ e $\mat{V}$ --- consultas, chaves e valores --- são projeções
lineares da entrada. Ao contrário da recorrência, o caminho entre duas posições
quaisquer tem comprimento constante, o que elimina o problema do gradiente evanescente
por outra via. O custo é quadrático em $T$.

Nenhum dos modelos treinados neste livro é um Transformer. Ele aparece, contudo, de
forma decisiva: o modelo de linguagem da \cref{parte:explicabilidade} é um Transformer
com centenas de bilhões de parâmetros, e compreender a \cref{eq:atencao} é pré-requisito
para entender o que ele faz e o que ele não faz. Essa discussão é o objeto do
\cref{cap:llm}.

\begin{destaque}[Prática]
O caderno \texttt{3\_introduction\_to\_unsupervised\_learning}, descrito no
\cref{cap:pratica}, treina um autoencoder LSTM sobre janelas do 3W e usa o erro de
reconstrução como escore de anomalia \citep{Lopes2026Intro3W}. Suas escolhas de
estabilidade numérica --- ativação \texttt{tanh}, inicialização ortogonal e recorte de
gradiente --- tornam concreto por que treinar recorrentes sobre sequências longas exige
cuidados que redes densas dispensam.
\end{destaque}

\resumocapitulo{%
As arquiteturas empregadas neste livro formam um repertório pequeno e bem delimitado.
Camadas densas fornecem os classificadores binários. Células LSTM, com seu estado de
célula propagado por soma, tornam viável o processamento de janelas de centenas de
amostras e sustentam os autoencoders recorrentes dos
\cref{cap:dual,cap:mundoaberto,cap:mahalanobis}. Convoluções unidimensionais, com pesos
compartilhados no tempo e campo receptivo controlado pela profundidade e pela
subamostragem, sustentam a U-Net de segmentação do \cref{cap:segmentacao}, cujas
conexões de atalho preservam a informação de localização temporal. Autoencoders
comprimem a entrada em uma representação latente e permitem detectar anomalias por erro
de reconstrução; suas variantes variacional e de agrupamento profundo não produziram
ganho mensurável, o que indica que o gargalo estava no objetivo de reconstrução e não na
arquitetura. O mecanismo de atenção, embora não treinado aqui, é a base dos modelos de
linguagem da Parte IV.}

\begin{exercicios}
\item Calcule o número de parâmetros de uma camada LSTM com $n_h$ unidades recebendo
      entrada de dimensão $n_x$. Aplique ao codificador do \cref{cap:mahalanobis}
      (LSTM(256) seguido de LSTM(128), entrada de 4 canais) e compare com o número de
      parâmetros de um MLP de mesma capacidade de saída sobre a entrada achatada de
      \num{800} dimensões.

\item Mostre, a partir da \cref{eq:lstm-cell}, que quando $\vect{f}_t = \vect{1}$ e
      $\vect{i}_t = \vect{0}$ para todo $t$ em um intervalo, o gradiente de
      $\vect{c}_{t_2}$ em relação a $\vect{c}_{t_1}$ é a identidade. Explique por que
      esse é o mecanismo que resolve o gradiente evanescente.

\item Calcule o campo receptivo de uma pilha de quatro convoluções unidimensionais com
      núcleo 3 e dilatação $1, 2, 4, 8$. Compare com quatro convoluções de núcleo 3 sem
      dilatação. Discuta quando convoluções dilatadas seriam preferíveis à subamostragem
      da U-Net.

\item Um autoencoder com $d = T \cdot S$ pode aprender a função identidade e ter erro de
      reconstrução nulo para qualquer entrada, inclusive anômala. Discuta três formas de
      impedir isso além da redução de $d$, e relacione cada uma a um mecanismo de
      regularização apresentado neste capítulo.

\item Derive a \cref{eq:vae} a partir do limite inferior da evidência (ELBO). Explique
      qual o papel do coeficiente $\beta$ e o que acontece nos limites $\beta \to 0$ e
      $\beta \to \infty$.

\item Implemente uma U-Net unidimensional com dois níveis e treine-a para segmentar
      regiões anômalas em uma instância sintética. Em seguida, remova as conexões de
      atalho, retreine e compare o IOU. Quantifique a degradação.

\item A complexidade da atenção é $O(T^2 d)$. Para $T = \num{86400}$ (um dia a
      \SI{1}{\hertz}) e $d = 128$, estime o número de operações de uma única camada de
      atenção. Discuta por que isso inviabiliza a aplicação direta de Transformers a
      séries temporais brutas de poço e cite duas estratégias da literatura para
      contornar o problema.
\end{exercicios}

\begin{leituras}
\item \citet{Hochreiter1997} --- artigo original da LSTM. Denso, mas a
      \cref{eq:lstm-cell} fica muito mais clara depois de lê-lo.
\item \citet{SCHMIDHUBER201585} --- panorama histórico do aprendizado profundo, útil
      para situar cada arquitetura em sua motivação original.
\item \citet{Ronneberger2015} --- a U-Net. A figura em forma de U do artigo é uma das
      mais reproduzidas da área, e por bons motivos.
\item \citet{Vaswani2017} --- o Transformer.
\item \citet{Kingma2014} e \citet{VAE} --- autoencoders variacionais.
\item \citet{srivastava14a} --- \ing{dropout}, com a análise que justifica sua
      interpretação como \ing{ensemble} implícito.
\end{leituras}
